\scanpage{4}
\uline{Tapis Cartier}

1)
\diagram{\begin{tikzcd}[ampersand replacement=\&]
 A \arrow[r,shift left,"f"] \& B \arrow[l,shift left,"g"]
\end{tikzcd}}
$g$ adjoint : \deleted{\unsure{gauche}}droite de $f$, d'où morphismes d'adjonction
\[
 \mathrm{id}_A\xrightarrow{\beta}gf,\qquad fg\mathrel{\overset{\alpha}{\unsure{-|\to}}}\mathrm{id}_B \tag*{(1.1)}
\]
On pose
\[
 \varphi=fg:\ B\mathrel{\unsure{-|\to}} B \tag*{(1.2)}
\]
On a donc
\[
 \boxed{\varphi\xrightarrow{\alpha}\mathrm{id}_B,\quad\varphi\xrightarrow{\lambda}\varphi^2} \tag*{(1.3)}
\]
où
\[
 \lambda\deleted{\unsure{=}}:fg\longrightarrow fgfg\qquad \lambda=f*\beta*g \tag*{(1.4)}
\]
\deleted{Alors} $\varphi$ satisfait \uline{conditions de Godement}

\diagram{\begin{tikzcd}[ampersand replacement=\&,column sep=small]
 \text{(1.5) (\uline{unitarité})} \& \varphi \arrow[r,"\lambda"] \arrow[rr,bend right=30,"\mathrm{id}_{\varphi}"'] \&
 \varphi^2 \arrow[r,shift left,"\varphi*\alpha\deleted{\unsure{*\varphi}}"] \arrow[r,shift right,"\alpha*\varphi"'] \& \varphi
\end{tikzcd}}
commutatif : $(\varphi*\alpha)\circ\lambda=\mathrm{id}_{\varphi}$

\diagram{\begin{tikzcd}[ampersand replacement=\&,column sep=small]
 \text{(1.6) (\uline{associativité})} \& \varphi\arrow[r,"\lambda"] \& \varphi^2
 \arrow[r,shift left,"\varphi*\lambda"] \arrow[r,shift right,"\lambda*\varphi"'] \& \varphi^3
\end{tikzcd}}
commutatif : $(\varphi*\lambda)\circ\lambda=(\lambda*\varphi)\circ\lambda$

On définit la catégorie
\[
 K=K(\varphi,\alpha,\lambda) \tag*{(1.7)}
\]
comme formée des couples
\[
 (X,u)\qquad X\in\mathrm{Ob}\,B,\quad \boxed{u:X\mathrel{\unsure{-|\to}}\varphi(X)} \tag*{(1.8)}
\]
satisfaisant les conditions de commutation
\diagram{\resizebox{\linewidth}{!}{\(\left\{\begin{tikzcd}[ampersand replacement=\&,column sep=small,row sep=large]
 \text{(\deleted{\unsure{1.9}}1.10)} \& {}\arrow[d,phantom,"{\left\{\vphantom{\rule{0pt}{2cm}}\right.}" description] \& X\arrow[r,"u"] \& \varphi(\deleted{\unsure{Y}}X)
 \arrow[r,shift left,"\varphi(u)"] \arrow[r,shift right,"\lambda(X)"'] \& \varphi^2(X)
 \& \lambda\circ u=\varphi(u)\circ u \quad\text{``associativité''}\\
 \text{(\deleted{\unsure{1.10}}1.9)} \& {} \& X\arrow[r,"u"] \& \varphi(X)\arrow[r,"\alpha(X)"] \& X
 \& \alpha(X)\circ u=\mathrm{id}_X\quad\text{``unitarité''}
\end{tikzcd}\right.\)}}
On obtient un foncteur naturel
\[
 \boxed{h:A\longrightarrow K(\varphi,\alpha,\lambda)} \tag*{(1.11)}
\]
par
\[
 Y\longmapsto \bigl(\deleted{\unsure{f}}f(Y),\quad f(Y)\xrightarrow{f*\beta}\varphi f(Y)=fgf(Y)\bigr).
\]
\margin{vérifier compatibilité (1.9)}
% archive sheet number: 3
